% Part: first-order-logic
% Chapter: beyond
% Section: introduction

\documentclass[../../../include/open-logic-section]{subfiles}

\begin{document}

\olfileid{fol}{byd}{int}

\olsection{جائزہ}

پہلے درجے دی منطق دلچسپی دا واحد منطقی نظام نہیں:
ایہدیاں کئی توسیعاں تے وکھرے روپ نیں۔ کسے منطق وچ
عام طور اُتے زبان دی رسمی وضاحت ہوندی اے؛ عموماً،
پر ہمیشہ نہیں، اک استنباطی نظام ہوندا اے؛ تے عموماً،
پر ہمیشہ نہیں، اک مقصود معنویات ہوندی اے۔ پر اصطلاح
”منطق“ دی ایس فنی ورتوں نال اک واضح سوال اٹھدا اے:
جیہڑیاں منطقاں پہلے درجے دی منطق نہیں، انہاں دا
”منطق“ لفظ نال کیہڑا رشتہ اے جدوں ایہ لفظ عام فہم
یا فلسفیانہ معنی وچ ورتیا جاوے؟ ہیٹھاں بیان کیتے سارے
نظام کسے نہ کسے طرح دی استدلالی سوچ نوں نمونہ بنان
لئی تیار کیتے گئے نیں؛ کی اسیں دس سکدے آں کہ انہاں
نوں منطقی کیہڑی گل بناؤندی اے؟

ایہدے آسان جواب سامنے نہیں آندے۔ ”منطق“ لفظ وکھ وکھ
طریقیاں تے سیاقاں وچ ورتیا جاندا اے، تے ”سچائی“ وانگ
ایس تصور دا وی کئی فلسفیانہ نقطۂ نظراں توں تجزیہ
ہویا اے۔ مثال لئی کوئی منطقی استدلال دا مقصد ایہ طے
کرنا سمجھ سکدا اے کہ کیہڑے بیان لازماً سچے نیں، تجربے
توں آزاد سچے نیں، غیر منطقی اصطلاحاں دی تشریح توں
آزاد سچے نیں، اپنی صورت دی وجہ توں سچے نیں، یا
لسانی رواج دی وجہ توں سچے نیں؛ تے انہاں ہر اک
تصور نوں ہور خاصی وضاحت دی لوڑ اے۔ جے توجہ صرف
ریاضی وچ ورتی جاندی منطق تک وی محدود رکھیے، تاں
وی ایہدی حد بارے بہت گھٹ اتفاق اے۔ مثال لئی
\textit{Principia Mathematica} وچ رسل تے وائٹ ہیڈ
نے فریگے توں شروع ہوئی {\em منطقیت} دی روایت وچ
ریاضی نوں منطق دی بنیاد اُتے تشکیل دین دی کوشش
کیتی۔ انہاں دا منطقی نظام اعلیٰ قسم دی منطق دا
اک روپ سی، جیہڑا ہیٹھاں بیان کیتے نظام نال ملدا
جلدا اے۔ آخرکار انہاں نوں ایداں دے مسلّمے شامل
کرن پئے جیہڑے اکثر معیاراں مطابق خالص منطقی نہیں
لگدے (خاص طور اُتے لامتناہیت دا مسلّمہ تے
تقلیل پذیری دا مسلّمہ)؛ پر فیر وی کوئی ایہ موقف
رکھ سکدا اے کہ اعلیٰ درجے دے استدلال دیاں کجھ
صورتاں منطقی منّیاں جانیاں چاہیدیاں نیں۔ ایس دے
برعکس، کوائن، جیہدی وجودیات ”قضیہاں“
(propositions) نوں گفتگو دیاں جائز شےواں نہیں
منّدی، دلیل دیندا اے کہ دوجے درجے دی تے
اعلیٰ درجے دی منطق اصل وچ سیٹ نظریے دیاں
ای صورتاں نیں، بس منطق دا بھیس پایا ہویا اے؛
دوجے لفظاں وچ، محمولاں اُتے مسوّریت والے نظام
خالص منطقی نہیں۔

فی الحال ایداں دے فلسفیانہ سوال کسے ہور ویلے
لئی چھڈ دینا بہتر اے، تے ہیٹھلے نظاماں نوں بس
استدلال دیاں وکھ وکھ قسماں دی رسمی مثالی
صورتاں سمجھنا اے، چاہے اوہ منطقی ہون یا نہ ہون۔

\end{document}
